\documentclass[11pt]{article}

\usepackage[T1]{fontenc}
\usepackage[utf8]{inputenc}
\usepackage{lmodern}
\usepackage{microtype}
\usepackage[a4paper,margin=1.08in]{geometry}
\usepackage{amsmath,amssymb,amsthm,mathtools,bm}
\usepackage{mathrsfs}
\usepackage{booktabs,array,longtable}
\usepackage{enumitem}
\usepackage{xcolor}
\usepackage{tikz}
\usepackage{pgfplots}
\usepackage{hyperref}
\usepackage{cleveref}
\usepackage{fancyhdr}
\usepackage{lastpage}

\pgfplotsset{compat=1.18}
\usetikzlibrary{arrows.meta,positioning,calc,decorations.pathreplacing}

\definecolor{navy}{RGB}{25,48,85}
\definecolor{teal}{RGB}{26,111,112}
\definecolor{rust}{RGB}{151,67,42}
\definecolor{lightgray}{RGB}{245,247,249}
\definecolor{midgray}{RGB}{90,98,108}

\hypersetup{
  colorlinks=true,
  linkcolor=navy,
  citecolor=teal,
  urlcolor=rust,
  pdftitle={A Squared-Complex Framework for Square-Prefix Mobius Sums},
  pdfauthor={Fred Viole, OVVO Labs}
}

\pagestyle{fancy}
\fancyhf{}
\fancyhead[L]{\small Squared-complex square-prefix M\"obius sums}
\fancyhead[R]{\small Fred Viole}
\fancyfoot[C]{\small \thepage\ of \pageref{LastPage}}
\setlength{\headheight}{14pt}

\setcounter{tocdepth}{2}
\setcounter{secnumdepth}{3}
\setlist[itemize]{leftmargin=1.5em,itemsep=0.25em,topsep=0.35em}
\setlist[enumerate]{leftmargin=1.7em,itemsep=0.3em,topsep=0.35em}
\allowdisplaybreaks
\raggedbottom
\emergencystretch=2em
\widowpenalty=10000
\clubpenalty=10000
\displaywidowpenalty=10000

\newtheorem{theorem}{Theorem}[section]
\newtheorem{proposition}[theorem]{Proposition}
\newtheorem{lemma}[theorem]{Lemma}
\newtheorem{corollary}[theorem]{Corollary}
\newtheorem{conjecture}[theorem]{Open theorem}
\theoremstyle{definition}
\newtheorem{definition}[theorem]{Definition}
\newtheorem{remark}[theorem]{Remark}
\newtheorem{example}[theorem]{Example}
\newtheorem{caution}[theorem]{Caution}

\crefname{conjecture}{open theorem}{open theorems}
\Crefname{conjecture}{Open theorem}{Open theorems}

\newcommand{\eps}{\varepsilon}
\newcommand{\Pplus}{P^{+}}
\newcommand{\1}{\mathbf{1}}
\newcommand{\vect}[1]{\bm{#1}}
\newcommand{\ip}[2]{\left\langle #1,#2\right\rangle}
\newcommand{\norm}[1]{\left\lVert #1\right\rVert}
\newcommand{\R}{\mathbb{R}}
\newcommand{\N}{\mathbb{N}}
\newcommand{\Z}{\mathbb{Z}}
\newcommand{\C}{\mathbb{C}}
\newcommand{\Repart}{\operatorname{Re}}
\newcommand{\Impart}{\operatorname{Im}}
\newcommand{\card}{\operatorname{card}}

\newcommand{\statusbox}[2]{%
\begin{center}
\fcolorbox{#1}{lightgray}{\parbox{0.91\textwidth}{\small\textbf{#2}}}
\end{center}}

\newcommand{\leanxref}[3]{%
  \par\smallskip
  {\small\noindent\textcolor{teal}{\textbf{#1:}}\par
  \noindent\href{#3}{\nolinkurl{#2}}\par}
  \smallskip
}

\title{\vspace{-1.5em}\textbf{A Squared-Complex Framework for Square-Prefix M\"obius Sums}\\[0.45em]
\large Exact arithmetic, canonical cofactor geometry, lifetime refinement, and the elementary RH bridge}
\author{Fred Viole\\OVVO Labs}
\date{August 2026}

\begin{document}
\maketitle

\begin{abstract}
This paper develops an exact geometric framework for square-prefix values of the Mertens function.  At the complete cutoff
\[
 X_n=(n+1)^2-1,
\]
a squarefree integer cannot contain two prime factors larger than $n+1$.  This gives the exact smooth--transport decomposition
\[
 S_n=M(X_n)=A_n-T_n.
\]
Each squarefree source $m>1$ has a unique canonical factorization
\[
 q_m=\Pplus(m),\qquad c_m=\frac{m}{q_m},
\]
and the squared-complex point
\[
 \left(\frac{c_m+q_m}{2}+i\frac{q_m-c_m}{2}\right)^2
 =m+i\frac{q_m^2-c_m^2}{2}
\]
records the source integer and its canonical factor imbalance simultaneously.  Fixed cofactors trace explicit parabolas, the square-prefix cutoff is vertical, and the square-root smoothness cutoff is a moving parabola whose intersection with the $c$-curve occurs exactly at $q=R$.

The largest-prime parametrization has a sharp endpoint support law that is useful throughout the paper.  If $m\le x$ and $\Pplus(m)>\lfloor x/2\rfloor$, then the canonical cofactor is $c_m=1$ and $m$ itself is prime.  Equivalently, every composite $m\le x$ satisfies
\[
 \Pplus(m)\le\lfloor x/2\rfloor.
\]
Thus the $c=1$ prime fibre is the entire extreme-largest-prime population, while every genuine product fibre has $c\ge2$ and largest prime at most $\lfloor x/2\rfloor$.

For a fixed height threshold $\Lambda>0$, low canonical height has uniformly bounded occupancy in each square block, giving the unconditional bound $|S_N^{\mathrm{low}}(\Lambda)|\ll_\Lambda N$.  A lifetime refinement splits the high population into currently active survivors and completed deaths.  The death increments are supported on a bounded-width integer window in the doubled height $|q^2-c^2|$; divisor bounds imply subpolynomial death increments and an RH-scale local-energy bound for the accumulated death mass.  Consequently the unresolved square-block cancellation may be placed entirely on the active survivor operator
\[
 Z_\Lambda(t)=-\sum_c\mu(c)\,\mathcal K_{\Lambda,t}(c),
\]
where $\mathcal K_{\Lambda,t}(c)$ counts canonical primes $q$ with $cq\le (t+1)^2-1$ and $2\Lambda t<|q^2-c^2|$.

The paper proves the exact arithmetic and geometric reductions, the low-height bound, and the death-process divisor-window estimate.  The remaining power-saving estimate for the signed survivor, equivalently the native high-sector square-prefix estimate, is stated explicitly as an open theorem.  The accompanying Lean development machine-checks the exact finite identities and elementary estimates.  No project axiom and no unconditional proof of the Riemann Hypothesis is claimed.
\end{abstract}

\statusbox{navy}{Status convention.  Results labelled theorem, proposition, lemma, or corollary are proved from the stated definitions or standard cited results.  Numerical values are finite computations.  Statements labelled ``Open theorem'' are the remaining analytic obligations and are not asserted as established.}
\newpage
\tableofcontents
\newpage
%======================================================================
